\section*{結言}
今回の8件では、世界モデルの実用化と公開モデルの選択肢拡大が中心的な流れとなった。Genie 4は操作性とマルチモーダルな表現を訴求し、Hy-Nomasは単一GPUで試せる構成を打ち出した。LFM3も小規模・多用途の設計を示しており、利用者は性能だけでなく、公開条件や実行環境を含めて技術を選ぶ局面にある。

一方、ETH反証の主張は、AIによる研究支援の可能性を印象づけると同時に、形式化された成果でも独立した専門家の検証が不可欠であることを浮き彫りにした。Cloudflareの事例は、推論インフラの高速化だけでなく、顧客データの隔離と検証の重要性を示している。

今後は世界モデルの利用条件と実測性能、LFM3の実機評価、ETHに関する第三者検証の結果が注目点となる。速報的な発表、コミュニティの解釈、検証済みの成果を区別して追うことが重要である。
